\FloatBarrier\section{Tecnologias Utilizadas}
\paragraph{Empacotamento e execução:} A imagem da aplicação reúne Python, Flask, Gunicorn e o cliente PostgreSQL. Flask implementa as rotas e constrói a resposta. Gunicorn recebe conexões e distribui o trabalho entre processos e threads dentro de cada réplica. Essa concorrência local é distinta da escala do Kubernetes, que acrescenta pods completos ao conjunto atendido pelo Service. O cliente psycopg abre a conexão usada para a operação SQL. Na implementação atual, a conexão é obtida por requisição, o que permite uma nova tentativa em consulta posterior à recuperação do banco.

\paragraph{Orquestração e configuração:} Minikube fornece o ambiente local em que é possível observar o ciclo de vida dos pods. O Deployment declara a quantidade desejada de réplicas da API e usa um ReplicaSet para mantê-la. O StatefulSet atende a outra necessidade: associar a instância do PostgreSQL a uma identidade e a uma solicitação de volume estáveis. O ConfigMap distribui parâmetros compartilhados, enquanto o Secret fornece os valores de autenticação aos consumidores. A Downward API entrega os dados individuais de cada pod. Essa combinação permite que instâncias originadas da mesma imagem se identifiquem de forma distinta ao cliente.

\paragraph{Saúde, isolamento e elasticidade:} BusyBox executa a espera inicial por TCP e também fornece ferramentas leves aos pods de teste. O kubelet consulta sondas distintas para decidir se o processo deve receber tráfego ou ser reiniciado. No PostgreSQL, as sondas usam a ferramenta de prontidão do próprio banco. O Calico aplica as regras de rede entre pods. Para a escala automática, o metrics-server fornece medidas de recursos e o controlador HPA ajusta o tamanho do Deployment de acordo com a utilização média de CPU. O HPA é opcional: mantê-lo separado da implantação base permite atribuir a expansão manual à ação do operador e a expansão sob carga à decisão do controlador.
\paragraph{Unidades de execução:} Dentro de cada pod da API, o contêiner BusyBox termina sua verificação antes que o contêiner da aplicação inicie. Só este último atende ao navegador. Já o pod temporário de carga permanece ativo durante o ensaio do HPA e gera as requisições que provocam consumo de CPU. A separação permite encerrar a carga sem reiniciar os servidores. Os clientes de teste da política de rede têm outra finalidade: variar o rótulo de origem enquanto se mantém o destino PostgreSQL.

\paragraph{Coordenação entre controladores:} O Deployment ajusta o número de réplicas por meio do ReplicaSet, enquanto o Service encaminha conexões aos pods aptos. O StatefulSet preserva a associação entre a identidade do banco e o PVC durante a recriação.
